%% ============================================================
%%  §4  The LLBC Deep Embedding
%% ============================================================

\section{The LLBC Deep Embedding}
\label{sec:llbc}

This section presents the deep embedding of LLBC into Lean~4. The embedding
consists of three parts: the AST types (\S\ref{sec:llbc-ast}), the interpreter
(\S\ref{sec:interpreter}), and the pipeline from Rust source to a verified
theorem (\S\ref{sec:pipeline}).

\subsection{The Verified Rust Fragment}
\label{sec:rust-fragment}

We verify a well-defined subset of safe Rust. The fragment is:

\begin{center}
\small
\begin{tabular}{@{}ll@{}}
\toprule
\textbf{Feature}             & \textbf{Status} \\
\midrule
Scalar integer types (I8/I16/I32/I64/U8/U16/U32/U64) & Supported \\
Boolean values               & Supported \\
Arithmetic with overflow checking & Supported \\
Comparisons, boolean logic   & Supported \\
Local variables (indexed)    & Supported \\
Assignments, sequencing      & Supported \\
\texttt{if}/\texttt{else}    & Supported \\
\texttt{loop}/\texttt{break}/\texttt{continue} & Supported \\
Function calls (cross-function via \texttt{FunEnv}) & Supported \\
\midrule
Struct fields, tuples        & Partial (projection in Charon; deref only in our AST) \\
\texttt{\&mut T} (mutable references) & Modelled via \texttt{FBPair} (\S\ref{sec:fbpair}) \\
\midrule
Raw pointers, \texttt{unsafe} blocks & \textbf{Not supported} \\
Closures, trait objects      & \textbf{Not supported} \\
Heap allocation (\texttt{Vec}, \texttt{Box}) & \textbf{Not supported} \\
Concurrency primitives       & \textbf{Not supported} \\
Lifetimes beyond \texttt{FBPair} & \textbf{Not supported} \\
\bottomrule
\end{tabular}
\end{center}

\noindent
This fragment matches what Charon can extract from safe Rust without
lifetimes or generic trait dispatch. Every function in the artifact
was checked to fall within this fragment before verification began.

\subsection{Trust Hierarchy}
\label{sec:tch}

The paper emphasises that the Lean kernel is the only trust anchor for
\emph{proof correctness}. However, the end-to-end pipeline introduces
additional trusted components for \emph{translation fidelity}:

\begin{center}
\small
\begin{tabular}{@{}lll@{}}
\toprule
\textbf{Component}    & \textbf{What it is trusted for} & \textbf{Verification status} \\
\midrule
Lean kernel           & All proof checks (258 theorems) & Formally verified \\
Mathlib               & Imported lemmas (\texttt{omega}, \texttt{Nat.fib}, etc.) & Community-reviewed \\
\midrule
Charon v0.1.197       & Rust MIR $\to$ LLBC JSON extraction & Unverified (external tool) \\
\texttt{charon2lean.py} & LLBC JSON $\to$ \texttt{LLBCFunDef} translation & Unverified (463 lines Python) \\
\bottomrule
\end{tabular}
\end{center}

\noindent
For the \textbf{hand-crafted} definitions in \texttt{ElabSpec.lean}
(F1--F16), the only components trusted are the Lean kernel and Mathlib:
the \texttt{LLBCFunDef} terms were written directly by the author and
inspected line by line against the Rust source. For \textbf{Charon-generated}
definitions (\texttt{CharonSpec.lean}, \texttt{NumIntegerSpec.lean}), the
trust chain extends to Charon and \texttt{charon2lean.py}: if either
mis-translates an instruction, the Lean kernel will verify a property about
an incorrect AST. This is analogous to the trusted code generator in Aeneas,
with the difference that our translator (Python, 463 lines) is much smaller
and auditable than a full compilation pipeline.

Mitigating this: (a)~every generated definition was manually reviewed against
the Charon JSON before proofs were written; (b)~ground-instance \rfl proofs
act as executable tests — if the AST is wrong, the kernel will not reduce to
the expected value.

\subsection{AST Types}
\label{sec:llbc-ast}

The LLBC type hierarchy in lean-pl-verify follows the structure of the Charon
intermediate representation~\cite{charon2024}. All types are defined as
inductive types in Lean~4.

\begin{lstlisting}[caption={LLBC type hierarchy.},label={lst:llbc-types}]
-- Scalar type tag (we support I32 and Bool in this fragment)
inductive LLBCTy where
  | I32 | Bool | Unit

-- A place is a local variable slot
inductive LLBCPlace where
  | var : Nat -> LLBCPlace

-- Operands are values that can be read from the current state
inductive LLBCOperand where
  | copy  : LLBCPlace -> LLBCOperand   -- borrow the value
  | move_ : LLBCPlace -> LLBCOperand   -- consume the value
  | const : LLBCLit   -> LLBCOperand   -- literal constant

-- Literals carry a type tag for scalar literals
inductive LLBCLit where
  | int  : Int -> LLBCTy -> LLBCLit
  | bool : Bool -> LLBCLit

-- Binary operators supported by the interpreter
inductive BinOp where
  | add | sub | mul | div | rem
  | eq | ne | lt | le | gt | ge
\end{lstlisting}

\texttt{LLBCPlace} is intentionally minimal: in the LLBC fragment we target
(safe Rust without raw pointers), every storage location is a numbered local
variable. The type supports extension to projections (\eg struct fields, tuple
components) but we leave that for future work.

\begin{lstlisting}[caption={LLBC rvalue and statement types.},label={lst:llbc-stmt}]
inductive LLBCRValue where
  | use    : LLBCOperand -> LLBCRValue
  | unOp   : UnOp -> LLBCOperand -> LLBCRValue
  | binOp  : BinOp -> LLBCOperand -> LLBCOperand -> LLBCRValue

inductive LLBCStmt where
  | assign   : LLBCPlace -> LLBCRValue -> LLBCStmt -> LLBCStmt
  | ite      : LLBCOperand -> LLBCStmt -> LLBCStmt -> LLBCStmt
  | loop     : LLBCStmt -> LLBCStmt
  | break_   : LLBCStmt
  | return_  : LLBCStmt
  | seq      : LLBCStmt -> LLBCStmt -> LLBCStmt
  | skip     : LLBCStmt

structure LLBCFunDef where
  name    : String
  params  : List LLBCTy
  locals  : List LLBCTy
  retTy   : LLBCTy
  body    : LLBCStmt

abbrev FunEnv := List LLBCFunDef
\end{lstlisting}

The \texttt{assign dst rvalue cont} constructor makes the continuation explicit:
the statement following an assignment is the third constructor field rather than
a separate sequencing node. This representation, adopted from CPS-style IRs,
slightly reduces the syntactic weight of typical function encodings; it means
that a function body like ``\texttt{x = a+b; y = x*2; return}'' is written
as a single nested \texttt{assign} rather than a chain of \texttt{seq} nodes.
Both styles are used in the artifact; \texttt{seq} is retained for the
post-loop continuation and explicit sequencing.

The \texttt{loop body} constructor takes a single body statement; the body
is expected to execute \texttt{break\_} to exit or fall through to restart.

\subsection{The Fuel-Parameterised Interpreter}
\label{sec:interpreter}

Lean~4's type theory requires all definitions to be terminating. LLBC loops
are potentially non-terminating, so the interpreter is parameterised by a
\emph{fuel} counter that decrements at each recursion step.

\begin{definition}[Fuel semantics]
\texttt{evalStmtFuel env fuel stmt locals} evaluates statement \texttt{stmt}
in function environment \texttt{env} with at most \texttt{fuel} recursive
steps, starting from local variable list \texttt{locals}. It returns
\texttt{.ok (sig, locals')} on success, where \texttt{sig} is a control-flow
signal (\texttt{.next}, \texttt{.break\_}, \texttt{.return\_}), or
\texttt{.error e} on panic.
\end{definition}

\begin{lstlisting}[caption={Core interpreter cases (simplified).},label={lst:interp}]
def evalStmtFuel (env : FunEnv) :
    Nat -> LLBCStmt -> Locals -> Except PanicReason (Signal × Locals)
  | _, .skip, s         => .ok (.next, s)
  | _, .break_, s       => .ok (.break_, s)
  | _, .return_, s      => .ok (.return_, s)

  -- Assignment: evaluate rvalue, update locals, run continuation
  | n+1, .assign dst rv k, s => do
      let v <- evalRValuePure rv s
      let s' := writePlacePure dst v s
      evalStmtFuel env n k s'

  -- Conditional: evaluate condition, dispatch to branch
  | n+1, .ite cond tb fb, s =>
      match evalOperandPure cond s with
      | some (.bool_ true)  => evalStmtFuel env n tb s
      | some (.bool_ false) => evalStmtFuel env n fb s
      | _ => .error .explicitPanic

  -- Loop: run body; .next restarts, .break_ exits
  | n+1, .loop body, s =>
      match evalStmtFuel env n body s with
      | .ok (.next,   s') => evalStmtFuel env n (.loop body) s'
      | .ok (.break_, s') => .ok (.next, s')
      | other             => other

  -- Explicit sequencing
  | n+1, .seq s1 s2, s =>
      match evalStmtFuel env n s1 s with
      | .ok (.next, s') => evalStmtFuel env n s2 s'
      | other           => other

  -- Out of fuel: treating this as a panic (not a logic error;
  -- correctness theorems ensure sufficient fuel is always supplied)
  | 0, _, _ => .error .explicitPanic
\end{lstlisting}

Two design decisions in this listing deserve explanation.

\textbf{Assignment continuation.} The \texttt{assign} case passes $n$ (not $n+1$)
to the continuation. This means each assignment step consumes one unit of fuel.
An intuitive way to think about this: fuel counts the number of evaluation
steps remaining, and each \texttt{assign} is one step.

\textbf{Loop fuel consumption.} In the \texttt{loop} case, the body is evaluated
with $n$ fuel (not $n+1$), and the recursive loop call also uses $n$ fuel.
This means the loop decrements by~1 per \emph{outer} loop invocation, not per
body step. The practical consequence is that a loop executing $k$ body iterations,
where each body needs $b$ fuel, requires $k \cdot b + k$ total fuel at the
outermost call. In the loop-correct theorems (\S\ref{sec:rust-proofs}), we
express this as $k + c \le \texttt{fuel}$ where $c$ accounts for the overhead
per iteration.

The helper functions \texttt{evalRValuePure} and \texttt{evalBinOpPure} are
pure (no monadic effects) and defined by case analysis; they return
\texttt{some v} on success and \texttt{none} on error. Binary operations check
overflow for \texttt{.add}, \texttt{.sub}, \texttt{.mul}:

\begin{lstlisting}[caption={evalBinOpPure for arithmetic.},label={lst:binop}]
def evalBinOpPure : BinOp -> Value -> Value -> Option Value
  | .add, .int a, .int b =>
      if minI32 <= a + b /\ a + b <= maxI32 then some (.int (a + b)) else none
  | .sub, .int a, .int b =>
      if minI32 <= a - b /\ a - b <= maxI32 then some (.int (a - b)) else none
  | .lt,  .int a, .int b => some (.bool_ (decide (a < b) = true))
  | .ge,  .int a, .int b => some (.bool_ (decide (a >= b) = true))
  -- ... further cases
\end{lstlisting}

The use of \texttt{decide (a < b) = true} rather than \texttt{a < b} is
deliberate: it produces a \texttt{Bool} via Lean's decidability mechanism,
which is kernel-evaluable on concrete integer values. The \texttt{Bool}-to-\texttt{Prop}
coercion is then handled by the \texttt{decide\_eq\_true} and
\texttt{decide\_eq\_false} lemmas in the conditional branch proofs.

\subsection{Function Evaluation}

\begin{lstlisting}[caption={Top-level function evaluation.},label={lst:evalfun}]
def buildLocals (f : LLBCFunDef) (args : List Value) : Locals :=
  -- allocate one slot for return value (index 0),
  -- then arg slots (index 1..params.length),
  -- then local variable slots (index params.length+1..end)
  let retSlot := [Value.unit]
  let argSlots := args.zipWith (fun _ v => v) f.params
  let localSlots := f.locals.map (fun _ => Value.unit)
  retSlot ++ argSlots ++ localSlots

def evalFunBody (env : FunEnv) (f : LLBCFunDef) (fuel : Nat)
    (locals : Locals) : Except PanicReason Value :=
  match evalStmtFuel env fuel f.body locals with
  | .ok (.return_, s) => (s[0]?).getD .unit |>.ok
  | .ok _             => .error .explicitPanic
  | .error e          => .error e

def evalFun (env : FunEnv) (f : LLBCFunDef) (fuel : Nat)
    (args : List Value) : RustM Unit Value :=
  fun _ =>
    match evalFunBody env f fuel (buildLocals f args) with
    | .ok v  => .ok (v, ())
    | .error e => .error e
\end{lstlisting}

The \texttt{at ()} syntax used in theorem statements (\eg
\texttt{evalFun [] f 30 args at ()}) is notation for applying the
\texttt{RustM Unit Value} computation to the trivial state \texttt{()}.
The state parameter is \texttt{Unit} because top-level function calls do not
share state with the caller; all state is local.

\subsection{Cross-Function Calls}
\label{sec:cross-calls}

The \texttt{FunEnv} list models a function environment that supports mutual
or cross-function calls. When a call instruction is evaluated, the interpreter
looks up the callee by name in \texttt{env} using \texttt{List.find?}:

\begin{lstlisting}[caption={Function call evaluation.},label={lst:funcall}]
  | n+1, .call fname args dest k, s =>
      match env.find? (fun f => f.name == fname) with
      | none   => .error .explicitPanic
      | some f =>
          let argVals := args.filterMap (fun op => evalOperandPure op s)
          match evalFunBody env f n argVals with
          | .ok v  => evalStmtFuel env n k (writePlacePure dest v s)
          | .error e => .error e
\end{lstlisting}

The cross-call uses $n$ fuel for the callee body, meaning the fuel is shared
between caller and callee steps. In the \texttt{square} $\to$ \texttt{mul}
example (\S\ref{sec:rust-proofs}), the concrete evaluation
\texttt{evalFun [mulFun] squareFun 20 [.int (-4)]} reduces by the
Lean kernel to \texttt{.ok (.int 16, ())} and the theorem is proved by \rfl.

\subsection{The Verification Pipeline}
\label{sec:pipeline}

The path from a Rust function to a Lean kernel-verified theorem proceeds
through five stages:

\begin{enumerate}[leftmargin=1.5em]
  \item \textbf{Write Rust source.} Functions are written in
        \texttt{examples/verified\_fns.rs}. Only safe Rust is used; the
        functions are simple enough that Charon can extract them without
        modifications.

  \item \textbf{Extract LLBC.} Run Charon on the Rust source to obtain the
        LLBC JSON.  The 463-line Python script \texttt{charon2lean.py}
        translates this JSON to \texttt{LLBCFunDef} Lean terms automatically
        (\S\ref{sec:eval-charon}).  Hand-crafted definitions in
        \texttt{ElabSpec.lean} (F1--F16) pre-date the pipeline and serve as
        the first verification layer; the pipeline is demonstrated on 18
        functions in \texttt{CharonSpec.lean}.

  \item \textbf{Encode as LLBCFunDef.} The LLBC AST is written as a Lean~4
        value of type \texttt{LLBCFunDef}. For example:

\begin{lstlisting}[caption={return42 function definition.},label={lst:r42}]
def return42Fun : LLBCFunDef := {
  name   := "return42"
  params := []
  locals := [.I32]
  retTy  := .I32
  body   := .assign (.var 0) (.use (.const (.int 42 .I32)))
              .return_
}
\end{lstlisting}

  \item \textbf{State the theorem.} The theorem asserts that \texttt{evalFun}
        on this definition satisfies the relevant \ProgramSpec. For
        \texttt{return42}:
\begin{lstlisting}
theorem elab_return42 :
    evalFun [] return42Fun 5 [] at () |=
      .pureOutput (. = .int 42) := by
  exact (.int 42, (), rfl, rfl)
\end{lstlisting}

  \item \textbf{Prove the theorem.} For pure functions, \rfl or
        $\langle\cdot, (), \texttt{rfl}, \texttt{rfl}\rangle$ closes the goal.
        For conditional functions, \texttt{simp only} with the relevant
        \texttt{decide\_eq\_true/false} lemma is used. For loops, the
        inductive invariant pattern described in \S\ref{sec:rust-proofs} is
        applied.
\end{enumerate}

This pipeline is entirely within Lean~4 and uses no external trusted components
beyond the Lean kernel and Mathlib. The 16 hand-crafted functions verified
(F1--F16) — \texttt{return42}, \texttt{id}, \texttt{neg}, \texttt{add},
\texttt{sub}, \texttt{max}, \texttt{abs}, \texttt{isZero}, \texttt{notGate},
\texttt{clamp}, \texttt{sumTo}, \texttt{fact}, \texttt{mul}, \texttt{min},
\texttt{square}, and \texttt{fib} — span the main categories of proof technique:
pure reduction (\rfl), overflow (\texttt{if\_pos}), conditional
(\texttt{decide\_eq\_true}), loop invariant (induction), and cross-function
(\rfl via kernel compute).
